\documentclass[11pt]{article}
\usepackage[margin=1in]{geometry}
\usepackage{amsmath,amssymb,amsthm}
\usepackage[round]{natbib}
\usepackage{booktabs}
\usepackage{graphicx}
\usepackage[hidelinks]{hyperref}
\usepackage{lmodern}
\usepackage[T1]{fontenc}
\usepackage{microtype}

\newtheorem{theorem}{Theorem}
\newtheorem{lemma}{Lemma}
\newtheorem{corollary}{Corollary}
\newtheorem{proposition}{Proposition}
\newtheorem{definition}{Definition}
\newtheorem{remark}{Remark}

\title{When Verifiers Backfire: A Cost--Correctness Theory of
Generation--Verification--Repair Loops for LLMs}

\author{Guidong Chen\thanks{Independent researcher. Email:
\href{mailto:yilingzitianlu@qq.com}{yilingzitianlu@qq.com}}}

\date{October 2026}

\begin{document}
\maketitle

\begin{abstract}
Test-time verification is rapidly becoming a standard component of
LLM agent systems: a generator produces an artifact, a verifier
accepts or rejects it, and rejections trigger repair.  Existing
theory (Dorner et al., ICLR 2026) characterizes how verifier
imperfection shapes pure resampling methods such as rejection
sampling and Best-of-$N$.  We extend this line of inquiry along three
axes that pure-resampling theory does not cover: (i) a
\emph{miss-asymmetry theorem} giving the closed-form exchange rate
$R = a/(1-b)$ between the two verifier error types when retries
redraw at the generator's rate---misses hurt
final precision $a/(1-b)$ times as much per unit as false rejections,
whose cost lands almost entirely in extra attempts; (ii) a
\emph{repair-degradation theorem} showing the exchange rate is
modulated by the retry success rate $q$: when $q$ falls below a
computable threshold $q^\ast$, false rejections become the dominant
harm and the asymmetry \emph{flips}; and (iii) an \emph{attribution
analysis} decomposing the terminal-error underprediction of pooled
per-attempt operating points into degraded retries (the dominant
channel) and error clustering (a residual).  We validate the
theory on HumanEval (164 problems, six arm configurations), pairing
every looped arm's terminal outcome against its own first attempt.
Oracle (unit-test) verification never loses a correct solution and
rescues $14$/$10$ wrong ones (McNemar $p < 10^{-3}$); LLM judges
with pooled operating points $(a,b) \approx (0.78\text{--}0.86,
0.74\text{--}0.86)$ do the reverse---net $-12$ correct solutions for
the isolated self-judge ($p = 0.004$, robust to counting only
delivered outcomes), $-8$ for the cross-capability judge ($p =
0.039$, convention-dependent), and a non-significant $-4$ for the
context-sharing judge---and the flip is predicted from measured
parameters alone ($q \approx 0.3$ vs.\ $q^\ast \approx 0.6$--$0.75$).
A budget sweep sharpens the practical rule: every judge arm's
single-shot delivered stream is purer than no verification, and the
first retry already erases the gain for the self-judge.  The
$\approx 3.4\times$ pooled-rate underprediction of terminal leak
(task-level bootstrap 95\% CI $[2.1, 4.7]$) is fully explained by
degraded retries: pooled rates mask the per-round trajectory, and
no residual requires error clustering.  A matched-pool ablation
\emph{suggests} that actionable feedback raises repair success
($0.33$ vs.\ $0.19$; small sample), and we report an honest
negative result on verifier-context
self-verification collapse.
Together, the results turn ``add a verifier'' from an article of
faith into a quantitative design decision.
\end{abstract}

\section{Introduction}\label{sec:intro}

Production LLM agent systems increasingly wrap generation in a
verification loop: an executor produces an artifact, a verifier
judges it, rejections trigger repair, and the loop terminates when
the verifier accepts or the budget runs out.  This
generation--verification--repair (GVR) pattern underlies executor/
verifier pipelines, LLM-as-judge filters, and test-driven agent
loops.  The engineering question it raises is sharp and often
ignored: \emph{when does adding a verifier actually improve the
final answer, and when does it make things worse?}

Two literatures border this question without answering it.  The
self-correction literature documents empirically that intrinsic
self-correction frequently fails or degrades performance
\citep{stechly2024selfverify,kamoi2024when,huang2024cannot}, while
sound external verification helps.  The test-time-scaling theory of
\citet{dorner2026rocnreroll} proves that instance-level accuracy of
rejection sampling and Best-of-$N$ is characterized by the geometry
of the verifier's ROC curve---but their analysis covers \emph{pure
resampling}: rejection carries no feedback, retries are independent
redraws, and no guidance is given on where along the ROC curve a
verifier should operate.

This paper contributes a compact theory of the GVR loop and a
$164 \times 6$ empirical study that stresses it.  Our contributions:

\begin{enumerate}
\item \textbf{Miss-asymmetry theorem (§\ref{sec:asym}).}  For a GVR
loop with generator success $p$ and verifier operating point
$(a,b)$, the marginal damage of the miss rate $1-b$ to final
precision exceeds that of the false-rejection rate $1-a$ by the
factor $R = a/(1-b)$.  When retries redraw at the generator's rate,
false rejections are nearly free for correctness and are paid for
only in attempts ($R \approx 19$ at $(0.95,0.95)$): push catch rate
aggressively.  The advice is conditional---§\ref{sec:degrade} derives
the threshold $q^\ast$ below which it inverts.
\item \textbf{Repair-degradation theorem (§\ref{sec:degrade}).}  The
exchange rate above assumes retries redraw at the original success
rate.  Real repair attempts are degraded---conditioned on rejection,
retries draw from a much harder pool, with success $q \ll p$.  We
derive the modulated exchange rate
$R_q = R \cdot \frac{q(1-p)}{p(1-q)}$ and the flip threshold
$q^\ast = \frac{(1-b)p}{(1-p)a + (1-b)p}$ below which \emph{false
rejections become the dominant harm}.  Measured $q \approx 0.3$
against $q^\ast \approx 0.6$--$0.75$ predicts---from parameters
alone ($q$ itself estimated from a mixed rejection pool;
§\ref{sec:exp})---that the LLM-judge arms in our study backfire;
they do, most robustly the isolated self-judge (within-arm McNemar
$p = 0.004$).
\item \textbf{Pooled-rate attribution (§\ref{sec:cluster}).}
Pooled per-attempt operating points, the natural summary reported by
most evaluation studies, underpredict terminal leak by a consistent
$\approx 3.4\times$ across our judge arms (task-level bootstrap
95\% CI $[2.1, 4.7]$).  The cause is that pooled rates mask the
per-round success trajectory $p_k$: a degradation-only model at the
measured $p_k$ reproduces the gap in full, leaving no residual that
requires cross-instance miss clustering (the attribution split is
convention-dependent and not separably identified at this sample
size).
\item \textbf{Empirical validation and boundary conditions
(§\ref{sec:exp}).}  Oracle verification confirms the theory exactly
(against its own first attempts it never loses a correct solution
and rescues $14$/$10$ wrong ones at $+27\%$ cost); a matched-pool
ablation provides suggestive (small-sample) evidence that
actionable feedback raises repair success from $0.19$ to $0.33$;
and a
verifier-context experiment yields an honest negative
result: sharing the generator's reasoning with a same-model judge
did \emph{not} collapse its accuracy, delimiting the scope of
self-verification-collapse claims to self-correction loops rather
than verification contexts.
\end{enumerate}

\section{Related Work}\label{sec:related}

\paragraph{Self-correction and its critiques.}
Self-Refine \citep{madaan2023selfrefine} and Reflexion
\citep{shinn2023reflexion} popularized iterative self-improvement.
\citet{stechly2024selfverify} showed systematic collapse of
self-critique on formal tasks, with LLM verifiers exhibiting very
high false-rejection rates; \citet{huang2024cannot} and the critical
survey of \citet{kamoi2024when} reached compatible conclusions.  Our
theory gives these observations a probabilistic form: collapse
corresponds to failure-mode correlation between generator and
verifier (§\ref{sec:mixture}), and our experiment delimits where
that correlation does and does not appear.

\paragraph{Test-time scaling theory.}
\citet{dorner2026rocnreroll} prove that rejection sampling and
Best-of-$N$ accuracy is determined by verifier ROC geometry, and
that RS dominates BoN at fixed compute.  Their setting is pure
resampling: no feedback, no repair, no generator--verifier
correlation analysis, and no operating-point selection guidance.
Our paper is a deliberate three-axis extension: repair feedback
(§\ref{sec:repair}), operating-point calibration
(§\ref{sec:asym}), and degraded/correlated retries
(§\ref{sec:degrade},~\ref{sec:cluster}).  Related theory includes
inference-time alignment coverage analysis
\citep{huang2025bestofn}.  In the training-time regime,
\citet{plesner2026imperfect} study robustness of RLVR to reward
noise and explicitly leave asymmetric (FN/FP) noise to future work;
our Theorem~\ref{thm:asym} supplies the inference-time counterpart
of that asymmetry in closed form.

\paragraph{Verification techniques.}
Self-consistency \citep{wang2023selfconsistency} uses agreement as
the verification signal; process supervision
\citep{lightman2024verify} trains step-level verifiers;
LLM-as-judge evaluation is surveyed by \citet{zheng2023judging}.
Control-theoretic formalizations of self-correction dynamics appear
in \citet{wu2026cybercorrect}; they manage loop stability, whereas
we manage the verifier's error budget.  All of these techniques
select a point on the verifier's quality trade-off; our results say
\emph{which} point to select and when the loop tolerates a harsh
verifier.

\section{Theory}\label{sec:theory}

\subsection{Setup: the GVR loop}\label{sec:setup}

One \emph{attempt}: a generator produces an artifact; a verifier
either \emph{accepts} or \emph{rejects} it (with feedback); a
rejection triggers a new attempt, up to budget $N$.

\begin{definition}[Per-attempt rates]\label{def:rates}
Let correctness of an attempt be binary, with generator success
$p = \Pr(C)$.  The verifier has acceptance rate
$a = \Pr(\text{accept} \mid C)$ and catch rate
$b = \Pr(\text{reject} \mid \neg C)$.  Its two error types are the
\emph{miss} $1-b$ (accepting wrong work) and the \emph{false
rejection} $1-a$ (rejecting correct work).
\end{definition}

\begin{definition}[Attempt outcomes]\label{def:outcomes}
Per attempt: $s = pa$ (accept-correct),
$\ell = (1-p)(1-b)$ (\emph{leak}: accept-wrong), and
$r = 1 - s - \ell$ (reject).
\end{definition}

Terminal states: \emph{accepted} (split into correct and leak) and
\emph{unresolved} (budget exhausted).  We assume (A1) the verdict is
independent of the generation path given correctness and (A2)
attempts are i.i.d.\ with constant $(p,a,b)$; both are relaxed in
§\ref{sec:repair}--\ref{sec:cluster}.

\begin{lemma}[Single-pass composition]\label{lem:single}
For $N=1$, $\Pr(\text{deliver wrong}) = (1-p)(1-b)$: a wrong
deliverable requires the conjunction of two independent failures.
The false-rejection rate does not enter this quantity.
\end{lemma}

\begin{proof}
Direct computation of $\Pr(\neg C \wedge \text{accept})$.
\end{proof}

\begin{theorem}[Loop closed forms]\label{thm:loop}
Under (A1)--(A2), with budget $N \ge 1$:
\begin{align}
\Pr(\text{leak}) &= \ell\,\frac{1-r^N}{1-r}, \quad
\Pr(\text{accepted}) = (s+\ell)\frac{1-r^N}{1-r},\\
\Pr(\text{unresolved}) &= r^N, \qquad
\mathbb{E}[\text{attempts}] = \frac{1-r^N}{1-r}.
\end{align}
\end{theorem}

\begin{proof}
The loop is geometric: reaching attempt $k$ requires $k-1$
rejections (probability $r^{k-1}$); summing and telescoping gives
the claims.  The instance-level structure matches the ROC-based
characterization of \citet{dorner2026rocnreroll} for the pure
resampling case; we state it in outcome-decomposed form for use
below.
\end{proof}

\begin{theorem}[Budget--precision invariance]\label{thm:invariance}
For every $N\ge1$,
\[
\pi := \Pr(\text{wrong} \mid \text{accepted})
= \frac{(1-p)(1-b)}{pa + (1-p)(1-b)},
\]
independent of $N$: added budget converts unresolved mass into
accepted mass \emph{at fixed precision}.
\end{theorem}

\begin{proof}
Ratio of the leak and accepted probabilities in
Theorem~\ref{thm:loop}; the geometric factors cancel.
\end{proof}

\subsection{The miss-asymmetry theorem}\label{sec:asym}

Both verifier errors hurt $\pi$---misses first-order through $\ell$,
false rejections second-order: each forced retry buys a fresh chance
to leak.  Their marginal damages have a closed-form ratio.

\begin{theorem}[Miss-asymmetry]\label{thm:asym}
\[
\frac{\partial \pi / \partial (1-b)}{\partial \pi / \partial (1-a)}
= \frac{a}{1-b}.
\]
\end{theorem}

\begin{proof}
$\partial\pi/\partial(1-b) = (1-p)s/(s+\ell)^2$ and
$\partial\pi/\partial(1-a) = p\ell/(s+\ell)^2$ by the chain rule on
Definition~\ref{def:outcomes}; the ratio simplifies to $a/(1-b)$.
\end{proof}

At $(a,b)=(0.95,0.95)$, $R \approx 19$: one unit of catch rate is
worth $\sim$19 units of acceptance rate for precision.  Meanwhile
the cost channels are reversed: $1-a$ enters expected attempts
first-order, $1-b$ only through the negligible $\ell$ term.

\begin{corollary}[Calibration guideline]\label{cor:calibrate}
Under (A2), any verifier modification trading acceptance for catch
at an exchange rate better than $1{:}R$ strictly improves final
precision; its price is paid only in attempts.  Calibrate verifiers
aggressively toward catch rate; budget false rejections rather than
fear them.  The qualifier is essential: when retries are degraded
below $q^\ast$ (Theorem~\ref{thm:flip}) the effective exchange rate
drops below 1 and the advice inverts---protect acceptance.
\end{corollary}

\begin{corollary}[Majority voting]\label{cor:majority}
Replace a verifier by a majority vote of $k$ conditionally
independent verifiers (strict majority rejects; ties accept).  For
$a, b > 1/2$ the vote raises \emph{both} effective rates:
\[
b^{(k)} = \sum_{j>k/2}\binom{k}{j}b^j(1-b)^{k-j} > b, \qquad
a^{(k)} = 1 - \sum_{j>k/2}\binom{k}{j}(1-a)^j a^{k-j} > a,
\]
because majority rule makes false rejections \emph{rarer}, not
commoner---a lone objecting judge no longer suffices to reject.
Precision therefore improves through both channels, and generator
attempts do not increase (the rejection mass $r$ falls); the price
is $k$ verifier invocations per round: at the pooled self-judge
operating point of §\ref{sec:exp}, majority-of-three lowers
expected generator attempts from $1.38$ to $1.25$ while raising
expected verifier calls from $1.38$ to $3.75$---the cost moves from
generation to verification, and whether it is worth paying depends
on the two units' relative prices.  The catch-vs-acceptance
trade-off that $R$ prices belongs instead to unanimity-style rules
(reject if any judge objects), where $a^{(k)} = a^k$ falls sharply
while $b^{(k)} = 1-(1-b)^k$ rises.
\end{corollary}

\subsection{Correlated failures}\label{sec:mixture}

(A1) fails in a specific direction: a verifier that is the same
model, primed with the generator's context, tends to repeat the
generator's mistake.  Model this with a latent failure-mode mixture.

\begin{definition}[Failure-mode mixture]\label{def:mixture}
Conditioned on latent mode $M=m$ with prior $\pi_m$, generation
succeeds with $p_m$; the verifier accepts correct output with
$a_m$ and catches wrong output with $b_m$, independently within
mode.
\end{definition}

\begin{definition}[Effective rates]\label{def:beff}
With $\epsilon = \sum_m \pi_m(1-p_m)$, error mass
$\mu_m = \pi_m(1-p_m)/\epsilon$ and success mass
$\sigma_m = \pi_m p_m/(1-\epsilon)$:
$b_{\mathrm{eff}} = \sum_m \mu_m b_m$ is the error-mass-weighted
catch rate and $a_{\mathrm{eff}} = \sum_m \sigma_m a_m$ the
success-mass-weighted acceptance rate.
\end{definition}

\begin{theorem}[Decorrelation condition]\label{thm:decorr}
Attempt-1 outcomes obey the closed forms of §\ref{sec:setup} with
pooled rates $(p,a,b) \to (1-\epsilon,\, a_{\mathrm{eff}},\,
b_{\mathrm{eff}})$; in particular the accepted-stream error
\[
\pi_1 = \frac{\epsilon\,(1-b_{\mathrm{eff}})}
{\epsilon\,(1-b_{\mathrm{eff}}) + (1-\epsilon)\,a_{\mathrm{eff}}}
\]
improves on no verification (delivered error $\epsilon$) iff
\[
a_{\mathrm{eff}} + b_{\mathrm{eff}} > 1,
\]
i.e., iff the mass-weighted operating point lies above the chance
diagonal of ROC space: verification helps exactly when
$\Pr(\text{accept}\mid C) > \Pr(\text{accept}\mid \neg C)$.  Catch
alone does not suffice: if $b_{\mathrm{eff}} \to 0$ while
$a_{\mathrm{eff}} < 1$, verification is \emph{strictly harmful}
($\pi_1 = \epsilon/[\epsilon+(1-\epsilon)a_{\mathrm{eff}}] >
\epsilon$), degenerating to neutral theater only in the
rubber-stamp limit $a_{\mathrm{eff}} \to 1$.  Correlated failure has
an exact form: with $\nu = \sum_m \pi_m(1-b_m)$ the prior-weighted
miss rate,
\[
\Pr(\text{wrong} \wedge \text{miss})
= \epsilon(1-b_{\mathrm{eff}})
= \epsilon\nu + \mathrm{Cov}_m(1-p_m,\,1-b_m),
\]
the covariance over modes under prior $\pi_m$: decorrelation is
precisely the claim that this covariance be small.
\end{theorem}

For $N>1$ the mixture bites: with instance-persistent modes,
surviving to round $k$ conditions on the mode mix.  §\ref{sec:exp}
validates the attempt-1 form directly---every judge arm's pooled
operating point lies above the chance diagonal, and every arm's
single-shot delivered stream is purer than no verification---while
the terminal gap beyond degraded retries is the quantity the
covariance term would have to explain; our data bound it but do
not identify it.

\subsection{Repair feedback breaks i.i.d.\ --- favorably}\label{sec:repair}

\begin{proposition}[Repair factor]\label{prop:delta}
Let rejected attempts be followed by a repaired attempt succeeding
with $p' = p + \delta(1-p)$, $\delta \in [0,1]$ ($\delta=0$:
rejection sampling; $\delta=1$: perfect repair).  Then all
formulas of §\ref{sec:asym} with constant $p$ are conservative
(actual leak is upper-bounded), and the invariance of
Theorem~\ref{thm:invariance} breaks in the favorable direction:
precision improves with budget.
\end{proposition}

\begin{proof}[Proof sketch]
Sample a uniform $u$ once per attempt slot and couple a repaired
draw (correct iff $u \le p'$) with a fresh one (correct iff
$u \le p$); since $p' \ge p$, the repaired draw is correct whenever
the fresh one is.  Every round $k \ge 2$ therefore has masses
$s' = p'a \ge s$ and $\ell' = (1-p')(1-b) \le \ell$, so its
accepted stream is at least as pure,
$\pi' = \ell'/(s'+\ell') \le \pi$, and every constant-$p$ formula
is an upper bound on realized leak.  The terminal conditional
error at budget $N$ is a weighted mean of $\pi$ (round 1) and
$\pi'$ (rounds $\ge 2$) with the $\pi'$-weight strictly increasing
in $N$; it therefore decreases toward $\pi'$---the favorable break
of Theorem~\ref{thm:invariance}'s invariance.
\end{proof}

\subsection{Repair degradation flips the asymmetry}\label{sec:degrade}

Real repairs do not redraw at rate $p$: the retry pool is
conditioned on rejection, dominated by hard instances, with success
$q \ll p$.

\begin{definition}[Retry rate]\label{def:q}
$q := \Pr(\text{attempt } k{+}1 \text{ correct} \mid \text{attempt }
k \text{ rejected})$, $k \ge 1$.
\end{definition}

A falsely rejected \emph{correct} solution enters a sub-loop with
per-round success $q$; conditional on resolution it ends wrong with
\[
\pi_q = \frac{(1-q)(1-b)}{qa + (1-q)(1-b)}.
\]

\begin{theorem}[Modulated exchange rate]\label{thm:flip}
At the population level, the marginal damage ratio becomes
\[
R_q \;=\; \frac{a}{1-b}\cdot\frac{q(1-p)}{p(1-q)}.
\]
The asymmetry \emph{flips} (false rejections dominate) iff
\[
q < q^\ast \;=\; \frac{(1-b)p}{(1-p)a + (1-b)p}.
\]
\end{theorem}

\begin{proof}
A unit $\varepsilon$ of extra miss rate moves $\varepsilon(1-p)$
mass of wrong first attempts from ``caught'' (sub-loop wrong with
probability $\pi_q$) to ``delivered wrong'' (probability 1):
damage $\varepsilon(1-p)(1-\pi_q)$.  A unit of extra false-rejection
rate moves $\varepsilon p$ mass of correct first attempts from
``delivered correct'' to the sub-loop: damage $\varepsilon p \pi_q$.
Substituting $\pi_q = \ell_q/(s_q+\ell_q)$ with $s_q = qa$,
$\ell_q = (1-q)(1-b)$ and taking the ratio yields $R_q$; setting
$R_q = 1$ and solving for $q$ gives $q^\ast$.  For $q=p$ the
expression recovers Theorem~\ref{thm:asym}.
\end{proof}

\begin{corollary}[Backfire prediction]\label{cor:backfire}
If a verifier has $1-a$ non-negligible and measured $q < q^\ast$,
adding it to the loop reduces terminal accuracy below the
no-verification baseline, despite any plausible catch rate.  The
conversion rate of correct solutions into wrong ones is
$(1-a)\,p\,\pi_q$.
\end{corollary}

\subsection{Clustering and pooled-rate failure}\label{sec:cluster}

The pooled $(\hat a, \hat b)$ reported by evaluation studies
summarize attempts across instances.  Two measured phenomena make
pooled rates poor predictors of terminal leak: (i) retry degradation
(§\ref{sec:degrade}) multiplies per-retry leak chance from
$(1-p)(1-b)$ to $(1-q)(1-b)$; (ii) misses cluster on particular
wrong-solution instances that recur across rounds, so
instance-level modes with low $b_m$ carry disproportionate mass.
Both effects push in the same direction; §\ref{sec:exp} quantifies a
consistent $\approx 3.4\times$ underprediction of leak by
pooled-rate theory and shows that (i) alone reproduces it, leaving
(ii) unidentified at this sample size.

\section{Experiments}\label{sec:exp}

\paragraph{Setup.}
HumanEval \citep{chen2021humaneval}, all 164 problems.  Generator:
qwen-plus (temperature 1.0, max 1024 tokens), prompted to plan then
emit a Python block.  Budget $N=5$.  Six arms: \textsf{none}
(deliver first attempt); \textsf{oracle\_full} (unit tests run in a
subprocess sandbox; feedback = failing-test traceback);
\textsf{oracle\_verdict} (tests; feedback = ``failed'' only);
\textsf{judge\_hetero} (qwen-max judge, isolated context);
\textsf{judge\_self\_iso} (qwen-plus judge, isolated context);
\textsf{judge\_self\_ctx} (qwen-plus judge shown the generator's own
reasoning).  Every attempt is additionally scored by the oracle
\emph{silently} for instrumentation, so per-arm operating points
$(\hat a, \hat b)$, per-round success $p_k$, and leak/terminal
outcomes are all measured directly.  All 984 runs completed; total
cost 1.64M tokens.

\begin{table}[t]
\centering\small
\caption{Main results ($n{=}164$ per arm; 95\% Wilson CIs).
$\pi_{\mathrm{pred}}$ from Theorem~\ref{thm:invariance} with
measured $\hat p_1,\hat a,\hat b$;
$\pi_{\mathrm{obs}}$ = error conditional on delivery.}
\label{tab:main}
\begin{tabular}{lccccccccc}
\toprule
Arm & $\hat p_1$ & $\hat a$ & $\hat b$ & Acc. & Leak &
Unres. & $\mathbb{E}[k]$ & $\pi_{\mathrm{pred}}$ &
$\pi_{\mathrm{obs}}$\\
\midrule
none            & .91 & --   & --   & .91 & .091 & .000 & 1.00 & --   & --\\
oracle\_full    & .89 & 1.00 & 1.00 & .98 & .000 & .024 & 1.27 & .000 & .000\\
oracle\_verdict & .90 & 1.00 & 1.00 & .96 & .000 & .037 & 1.29 & .000 & .000\\
judge\_hetero   & .90 & .78  & .77  & .85 & .098 & .049 & 1.52 & .030 & .103\\
judge\_self\_iso& .90 & .78  & .74  & .82 & .116 & .061 & 1.51 & .037 & .123\\
judge\_self\_ctx& .90 & .86  & .86  & .87 & .061 & .067 & 1.45 & .018 & .065\\
\bottomrule
\end{tabular}
\end{table}

\subsection{Oracle verification: theory confirmed exactly}
Unit-test verification lifts accuracy from .91 to .98 at $+27\%$
attempts ($\mathbb{E}[k]=1.27$), with leak collapsing from $9.1\%$
to exactly $0.0\%$, as Theorem~\ref{thm:loop} requires for $b=1$.
Within its own first attempts the loop never loses a correct
solution and rescues $14$/$10$ wrong ones
(Table~\ref{tab:within}).  This is the clean empirical instantiation
of Corollary~\ref{cor:calibrate}: at $a = 1$ its $q$-qualification
is vacuous (there are no false rejections to convert), and budget
spent on catch rate pays off at near-parity cost.

\subsection{LLM judges lose correct solutions, as the degradation
theorem predicts}
All three judge arms land \emph{below} the no-verification baseline
(.85/.82/.87 vs.\ .91) despite pooled operating points that look
respectable.  Cross-arm comparisons, however, are confounded by
first-draw sampling: each arm drew its first attempts independently,
and $\hat p_1$ ranges $.89$--$.91$ (\textsf{none} happened to draw
the best of the six).  The causal comparison pairs each arm's
terminal outcome against \emph{its own} first attempt on the same
task.

\begin{table}[t]
\centering\small
\caption{Within-arm pairing: terminal outcome vs.\ the same arm's
first attempt, by task (McNemar exact).  ``Lost'' = first attempt
correct, terminal not delivered-correct; ``rescued'' = first attempt
wrong, terminal correct.  $p$: unresolved counted as failure (the
paper's accuracy convention); $p_{\mathrm{del}}$: unresolved tasks
excluded.  Effect = terminal accuracy minus own first-attempt
accuracy.}
\label{tab:within}
\begin{tabular}{lcccccc}
\toprule
Arm & Lost & Rescued & Net & $p$ & $p_{\mathrm{del}}$ & Effect\\
\midrule
oracle\_full      & 0  & 14 & $+14$ & $.0001$ & $.0001$ & $+8.5$pp\\
oracle\_verdict   & 0  & 10 & $+10$ & $.002$  & $.002$  & $+6.1$pp\\
judge\_hetero     & 10 & 2  & $-8$  & $.039$  & $.29$   & $-4.9$pp\\
judge\_self\_iso  & 14 & 2  & $-12$ & $.004$  & $.013$  & $-7.3$pp\\
judge\_self\_ctx  & 10 & 6  & $-4$  & $.46$   & $1.00$  & $-2.4$pp\\
\bottomrule
\end{tabular}
\end{table}

Table~\ref{tab:within} is the paper's cleanest contrast.  Oracle
verification never loses a correct solution (a unit acceptance rate
never re-rolls correct work) and rescues $14$/$10$ wrong first
attempts.  The isolated self-judge does the reverse, significantly
under both conventions ($p = .004$ and $.013$)---the one backfire
robust to how unresolved tasks are counted.  The cross-capability
judge is significant only under the accuracy convention (it fails
a Bonferroni correction across the three judge arms and loses
significance when unresolved tasks are excluded); the
context-sharing judge's deficit is a null.  We accordingly claim a
significant backfire only for \textsf{judge\_self\_iso}.

The mechanism is measured directly: false rejections
send correct solutions into a retry sub-loop whose conditional
success is $q \approx 0.27$--$0.39$ (Table~\ref{tab:pk}), so
correct$\to$wrong conversion occurs at the rate
Corollary~\ref{cor:backfire} predicts, $(1-\hat a)\hat p_1
\hat\pi_q$ with $q = \hat p_2$: predicted 15.7/5.6/10.0 conversions
for \textsf{self\_iso}/\textsf{self\_ctx}/\textsf{hetero} vs.\
observed 12/6/6---the correct order and the correct arm ranking.
Predictions are expectations; with ${\sim}30$ expected events the
binomial noise is ${\pm}5$.  (Measured $\hat p_2$ pools rejections
of genuinely wrong and of correct attempts---55 vs.\ 39 rejection
events in \textsf{hetero}---so it approximates, rather than
identifies, the correct-solution sub-loop rate; cf.\ Limitations,
§\ref{sec:discussion}.)  The flip thresholds are
$q^\ast \approx 0.75/0.59/0.73$; measured $q$ lies far below all of
them, so the backfire is predicted from parameters alone, before
any end-to-end run.

\begin{table}[t]
\centering\small
\caption{Per-round conditional success $p_k$ (rows: arms;
$n\ge10$). $q := \hat p_2$ (second-round conditional success), the
quantity used in §\ref{sec:exp}; \textsf{judge\_self\_ctx}'s
$p_5 = .00$ rests on 11 rounds in a heavily depleted pool.}
\label{tab:pk}
\begin{tabular}{lccccc}
\toprule
Arm & $p_1$ & $p_2$ & $p_3$ & $p_4$ & $p_5$\\
\midrule
oracle\_full    & .89 ($n{=}164$) & .33 (18) & .25 (12) & -- & --\\
oracle\_verdict & .90 (164) & .19 (16) & .23 (13) & .10 (10) & --\\
judge\_hetero   & .90 (164) & .39 (36) & .42 (24) & .13 (15) & .45 (11)\\
judge\_self\_iso& .90 (164) & .27 (37) & .37 (19) & .35 (17) & .40 (10)\\
judge\_self\_ctx& .90 (164) & .30 (30) & .32 (19) & .29 (14) & .00 (11)\\
\bottomrule
\end{tabular}
\end{table}

\subsection{Single-shot verification helps; retries eat the gain}
\label{sec:budget}
Truncating every run at budget $N$ and recomputing
conditional-on-delivery error traces the loop's purity curve
(Table~\ref{tab:budget}).

\begin{table}[t]
\centering\small
\caption{Conditional-on-delivery error $\pi$ by budget $N$
(runs truncated at $N$ rounds).  The no-verification baseline
(deliver the first attempt unverified) has $\pi = .091$.}
\label{tab:budget}
\begin{tabular}{lccccc}
\toprule
Arm & $N{=}1$ & $N{=}2$ & $N{=}3$ & $N{=}4$ & $N{=}5$\\
\midrule
judge\_hetero    & .031 & .050 & .074 & .092 & .103\\
judge\_self\_iso & .024 & .103 & .109 & .123 & .123\\
judge\_self\_ctx & .022 & .055 & .067 & .065 & .065\\
\bottomrule
\end{tabular}
\end{table}

Two readings.  First, at $N=1$ every judge arm's delivered stream is
far purer than the no-verification baseline ($.022$--$.031$ vs.\
$.091$), exactly as Theorem~\ref{thm:decorr} requires: each arm's
pooled operating point lies above the chance diagonal ($\hat a +
\hat b = 1.52$--$1.72$), so single-shot verification helps.
Second, purity degrades with budget: retries draw at $q \ll p$, so
each additional round admits increasingly impure accepted mass.
Budget is thus \emph{not neutral}: Theorem~\ref{thm:invariance}'s
invariance holds only when retries redraw at the generator's
rate---Proposition~\ref{prop:delta} shows it breaking favorably
when repairs improve, the degradation of §\ref{sec:degrade}
unfavorably when they do not---and the sweep realizes the
unfavorable side, where added budget \emph{actively} worsens the
delivered stream.  For \textsf{judge\_self\_iso} the
first retry alone erases the entire single-shot gain ($\pi$:
$.024 \to .103$, crossing the baseline $.091$).  Budget, like
strictness, is a $q$-dependent decision (G2, §\ref{sec:discussion}).

\subsection{Pooled operating points underpredict leak by
$\approx3.4\times$---fully explained by degraded retries}
Using measured pooled $(\hat a,\hat b)$ and $\hat p_1$,
Theorem~\ref{thm:invariance} predicts conditional-on-delivery error
of 3.0\%/3.7\%/1.8\%; observed are 10.3\%/12.3\%/6.5\%---per-arm
ratios of 3.4/3.3/3.6, geometric mean $3.4$ (task-level bootstrap
95\% CI $[2.1, 4.7]$, far from 1) (Table~\ref{tab:main}, last two
columns).  What does this gap mean?  A degradation-only model---the
measured per-round trajectory $p_k$ (Table~\ref{tab:pk}) with
constant $(\hat a, \hat b)$, no instance-level clustering---reproduces
it in full: predicted $\pi$ of .106/.134/.068 vs.\ observed
.103/.123/.065, slightly \emph{over}predicting every arm.  Degraded
retries are thus not merely the dominant channel but a
\emph{sufficient} one: the data require no clustering channel.  The
precise attribution share, however, is a modeling convention rather
than a measurement: across the two conventions we implement
(retries at constant $q = \hat p_2$ vs.\ at the measured $p_k$;
per-round survival) it spans 98--128\%, and at this sample size
degradation and clustering are not separably identified; a
per-instance binomial overdispersion check
($\chi^2/\mathrm{df} = 1.3$--$2.2$ across judge arms; the same wrong
instances do leak at $k{=}1$ through $k{=}5$) is directionally
consistent with mild miss clustering but does not separate the
channels.  The leak composition confirms that degradation carries
the gap: of the 16 leaks in \textsf{judge\_hetero}, 6 arise from
falsely rejected \emph{correct} first attempts---the conversion
channel of Theorem~\ref{thm:flip}---and 10 from wrong first attempts
that survived repeated verification; in the two self-judge arms the
conversion channel is the majority (12/19 and 6/10).
Instance-level modeling, as in \citet{dorner2026rocnreroll}, remains
a necessity for terminal-error prediction; but the first-order
lesson is simpler---pooled rates fail because they mask the
per-round trajectory, so measure $q$.

\subsection{Actionable feedback repairs better (matched pools)}
\textsf{oracle\_full} vs.\ \textsf{oracle\_verdict} share the same
retry pool (first-attempt-wrong problems), isolating feedback
content; with $b=1$ the pool is exactly the wrong-first-attempt
tasks, so this comparison does not inherit the mixed-pool
approximation of the judge arms.  Traceback feedback raises
second-round repair success
from .19 to .33 ($+74\%$ relative; $3/16$ vs.\ $6/18$ events,
Fisher exact $p = 0.45$) and lowers unresolved rate (.024 vs.\
.037).  The sample is small, so we read this as \emph{suggestive}
evidence on a literature tension rather than an adjudication:
\citet{stechly2024selfverify} found feedback detail immaterial with
sound verifiers on formal tasks, while Reflexion-style systems
premise feedback-driven repair; our $\delta$-framework
(Proposition~\ref{prop:delta}) can accommodate both regimes, and
code repair appears to sit in the $\delta>0$ regime.

\subsection{An honest negative result on context-induced collapse}
Contrary to the collapse motivation of §\ref{sec:mixture}, the
judge \emph{shown the generator's own reasoning}
(\textsf{judge\_self\_ctx}) was slightly \emph{more} accurate at
both operating points than the isolated same-model judge
($\hat a = .86$ vs.\ $.78$; $\hat b = .86$ vs.\ $.74$) and delivered
better terminal accuracy (.87 vs.\ .82).  These gaps are
non-significant---the $\hat b$ Wilson intervals overlap
($[.76,.92]$ vs.\ $[.63,.83]$)---and the context arm's within-arm
deficit is the weakest of the three (net $-4$, $p = 0.46$; $p =
1.0$ counting only delivered outcomes), so
we claim no evidence of context-induced collapse but cannot exclude
a moderate harmful effect either.  In this setting,
context-sharing did not corrupt verification; collapse phenomena
appear confined to self-\emph{correction} loops
\citep{stechly2024selfverify}, not verifier contexts.  We report
this as a boundary condition on the theory: the correlation that
destroys verification loops must be measured
(Theorem~\ref{thm:decorr}), not assumed.

\section{Discussion}\label{sec:discussion}

\paragraph{Design guidelines.}
\textbf{G1: oracle-first.}  Wherever a deterministic check exists
(tests, compilers, schemas), attach it to the verifier: it removes
leak entirely at bounded cost.  \textbf{G2: strictness must be
earned by repairability.}  Harsh verification is safe when retries
are cheap and successful ($q$ high); when repairs are degraded,
Corollary~\ref{cor:calibrate} inverts---acceptance must be protected.
Measure $q$ (one rejected-batch pilot) before choosing a verifier's
operating point \emph{or its budget} (§\ref{sec:budget}).
\textbf{G3: never trust pooled operating points}
for terminal-error prediction; report instance-level effects.
\textbf{G4: engineer feedback actionability ($\delta$)}---it is the
cheapest lever on repair success.

\paragraph{Human analogy.}
The same structure appears in human peer review: independent review
layers compose multiplicatively (the ``Swiss cheese'' model of
accident analysis), reviewer strictness is affordable exactly when
authors can repair well, and nitpicking a correct solution has its
price paid in rework---sometimes turning a working change into a
broken one, our conversion mechanism in miniature.

\paragraph{Limitations.}
Single generator family (Qwen) and one benchmark (HumanEval, 164
problems); judge ``heterogeneity'' is cross-capability, not
cross-family; leak rates rest on 10--19 instances per judge arm.
The feedback ablation rests on 3 vs.\ 6 repair events
(Fisher $p = 0.45$) and is suggestive only (its retry pool, unlike
the judge arms', is cleanly identified); the retry rate $q$ for
the judge arms is measured from a mixed rejection pool and only
approximates the correct-solution sub-loop rate; the
context-collapse negative
result is one setting with overlapping confidence intervals.  The
theory's identifiability assumptions (within-mode independence) are
idealizations.  The artifact code supports replication and
extension.

\section{Conclusion}

Verification is not free reliability.  Its benefit is governed by a
computable exchange rate between the verifier's two error types; its
danger---demonstrated here empirically at the rates the theory
predicts---appears when degraded repairs let false rejections
convert correct work into wrong.  The practical summary fits in one
sentence: \emph{spend your verifier budget on catching errors, but
first measure how well your generator recovers from being told it
was wrong.}

\paragraph{Reproducibility.}
Experiment code, prompts, per-attempt logs (984 runs, 1.64M tokens),
and analysis scripts are released at
\url{https://github.com/dianxia/verifier-backfire}.

\begin{thebibliography}{20}

\bibitem[Chen et al.(2021)]{chen2021humaneval}
M.~Chen, J.~Tworek, H.~Jun, et al.
\newblock Evaluating large language models trained on code.
\newblock \emph{arXiv:2107.03374}, 2021.

\bibitem[Dorner et al.(2026)]{dorner2026rocnreroll}
F.~E. Dorner, Y.~Chen, A.~F. Cruz, F.~Yang.
\newblock {ROC}-n-reroll: How verifier imperfection affects
test-time scaling.
\newblock \emph{ICLR}, 2026. arXiv:2507.12399.

\bibitem[Huang et al.(2024)]{huang2024cannot}
J.~Huang, X.~Chen, S.~Mishra, H.~S. Zheng, A.~W. Yu, X.~Song,
D.~Zhou.
\newblock Large language models cannot self-correct reasoning yet.
\newblock \emph{ICLR}, 2024.

\bibitem[Huang et al.(2025)]{huang2025bestofn}
A.~Huang, et al.
\newblock Is best-of-$n$ the best of them? Coverage, scaling, and
optimality in inference-time alignment.
\newblock \emph{arXiv:2503.21878}, 2025.

\bibitem[Kamoi et al.(2024)]{kamoi2024when}
R.~Kamoi, Y.~Zhang, N.~Zhang, J.~Han, R.~Zhang.
\newblock When can {LLM}s actually correct their own mistakes? {A}
critical survey of self-correction of {LLM}s.
\newblock \emph{TACL}, 12:1417--1440, 2024.

\bibitem[Lightman et al.(2024)]{lightman2024verify}
H.~Lightman, V.~Kosaraju, Y.~Burda, et al.
\newblock Let's verify step by step.
\newblock \emph{ICLR}, 2024.

\bibitem[Madaan et al.(2023)]{madaan2023selfrefine}
A.~Madaan, N.~Tandon, P.~Gupta, et al.
\newblock Self-refine: Iterative refinement with self-feedback.
\newblock \emph{NeurIPS}, 2023.

\bibitem[Plesner et al.(2026)]{plesner2026imperfect}
A.~Plesner, F.~Guzm\'an, A.~Athalye.
\newblock An imperfect verifier is good enough: Learning with noisy
rewards.
\newblock \emph{arXiv:2604.07666}, 2026.

\bibitem[Shinn et al.(2023)]{shinn2023reflexion}
N.~Shinn, F.~Cassano, E.~Berman, A.~Gopinath, K.~Narasimhan,
S.~Yao.
\newblock Reflexion: Language agents with verbal reinforcement
learning.
\newblock \emph{NeurIPS}, 2023.

\bibitem[Stechly et al.(2025)]{stechly2024selfverify}
K.~Stechly, K.~Valmeekam, S.~Kambhampati.
\newblock On the self-verification limitations of large language
models on reasoning and planning tasks.
\newblock \emph{ICLR}, 2025. arXiv:2402.08115.

\bibitem[Wang et al.(2023)]{wang2023selfconsistency}
X.~Wang, J.~Wei, D.~Schuurmans, et al.
\newblock Self-consistency improves chain of thought reasoning in
language models.
\newblock \emph{ICLR}, 2023.

\bibitem[Wu et al.(2026)]{wu2026cybercorrect}
Y.~Wu, Y.~Liu, Y.~Shu.
\newblock {CyberCorrect}: A cybernetic framework for closed-loop
self-correction in large language models.
\newblock \emph{arXiv:2605.17305}, 2026.

\bibitem[Zheng et al.(2023)]{zheng2023judging}
L.~Zheng, W.~Chiang, Y.~Sheng, et al.
\newblock Judging {LLM}-as-a-judge with {MT}-Bench and Chatbot
Arena.
\newblock \emph{NeurIPS Datasets and Benchmarks}, 2023.

\end{thebibliography}

\end{document}